\chapter{Theoretical Background}
\label{chap:background}

\section{Network Intrusion Detection Systems}
\label{sec:bg-nids}

A Network Intrusion Detection System inspects packets crossing a
network and emits an alert when those packets match a model of
malicious activity. The discipline grew out of host-based intrusion
detection in the late 1980s and was formalised in the late 1990s when
Vern Paxson published Bro --- the system later renamed
Zeek~\cite{paxson1999bro} --- and Martin Roesch published the first
version of Snort~\cite{roesch1999snort}. Both systems are still
maintained today and still influence the architecture of every
serious open-source NIDS, including this one.

A NIDS sits passively in front of the traffic it watches. It does not
forward, drop or modify packets; it only observes. Inline detection,
where the system actively blocks traffic, falls under the related
acronym NIPS (Network Intrusion Prevention System). \nidsname{} is a
NIDS in the strict sense: it can produce alerts but it cannot stop a
packet from reaching its destination. The reasoning is partly
philosophical and partly practical. Blocking traffic at the inspection
point creates a single point of failure for the whole network, and
production deployments overwhelmingly prefer to run detection in a
side-channel (a SPAN port, a TAP, a virtual interface) and to put
prevention behind a separate firewall layer.

The taxonomy of detection strategies usually splits along two axes:
the \emph{paradigm} (misuse-based versus anomaly-based) and the
\emph{granularity} (packet-level versus flow-level)~\cite{garcia2009anomaly}.
The paradigm axis is the more important one. Misuse-based or
signature-based detection encodes the patterns of known attacks; the
engine matches each packet against a curated set of rules. Anomaly-based
detection encodes a model of normal traffic; the engine flags packets
that deviate from that model. The granularity axis is more of an
implementation choice: a packet-level engine inspects each frame in
isolation, while a flow-level engine aggregates packets sharing the
same 5-tuple into a bidirectional flow before scoring. \nidsname{}
operates packet-level because the design priority was real-time
responsiveness; flow features such as inter-arrival time are kept as
streaming state inside the ML detector.

\section{Signature-Based Detection}
\label{sec:bg-signature}

Signature-based detection is the older and conceptually simpler
approach. The engine holds a database of patterns; for each packet it
iterates the patterns and emits an alert whenever one matches. The
patterns themselves are typically expressed in a domain-specific
language. Snort's rule syntax has become an informal standard and is
shared (with extensions) by Suricata~\cite{suricata}. A typical Snort
rule looks like this:

\begin{verbatim}
alert tcp $EXTERNAL_NET any -> $HOME_NET 22
  (msg:"SSH brute force attempt";
   flags:S;
   threshold:type both, track by_src, count 10, seconds 30;
   classtype:attempted-recon; sid:1000001; rev:1;)
\end{verbatim}

The header (\verb|alert tcp ... -> ... 22|) constrains the protocol,
direction and ports. The options inside the parentheses constrain
payload contents, TCP flag combinations, packet sizes, and so on. The
\verb|threshold| option turns a single-packet pattern into a temporal
one --- it fires only when a certain count of matches arrives within a
given window from the same source.

Signature-based detection has two well-understood strengths. It is
\emph{precise} (false-positive rate is low when rules are well-written)
and it is \emph{auditable} (each alert carries the identifier of the
rule that produced it, so an analyst can trace the decision back to a
human-written predicate). Its corresponding weakness is that it cannot
detect attacks for which no rule has yet been written. Every novel
attack technique has a window between its first appearance in the wild
and the day a signature for it ships in the rule
distribution~\cite{garcia2009anomaly}; during that window, a
signature-only engine sees nothing.

\nidsname{} uses a JSON-based match specification instead of the
Snort DSL. The decision was driven by two factors. First, a JSON match
maps cleanly to a Pydantic model on the backend and to a JavaScript
object on the frontend, which makes the rule-editing experience much
simpler than parsing and pretty-printing Snort syntax. Second, the
project does not aim to compete with Snort on rule expressiveness; it
aims to cover the categories of attack the validation scenario uses
(port scans, brute-force attempts, simple floods, content-based
exploits), and the JSON specification is sufficient for all of them.
Section~\ref{sec:design-signature} describes the exact schema.

\section{Anomaly-Based and Machine-Learning Detection}
\label{sec:bg-ml}

Anomaly-based detection inverts the signature paradigm. Instead of
asking ``does this packet look like a known attack?'' it asks ``does
this packet look like normal traffic?''. The engine learns a model of
the normal distribution from observed data and flags samples that fall
outside it. The literature on the topic is enormous; Buczak and Guven
counted over a hundred distinct proposals in a single survey
paper~\cite{buczak2016survey}, and Khraisat et~al.\ extended the count
in a more recent review~\cite{khraisat2019survey}.

I selected two algorithms from this body of work for \nidsname{}:
Random Forest and Isolation Forest. Both are tree-based ensemble
methods from scikit-learn~\cite{pedregosa2011scikit}, and both are
fast enough to run in the per-packet inference path without affecting
the overall throughput. They were chosen for their complementary
semantics rather than for any claim of state-of-the-art accuracy.

\subsection{Random Forest}
\label{sec:bg-rf}

The first algorithm came from Leo Breiman's 2001
paper~\cite{breiman2001random}. The construction is
straightforward to describe but surprisingly hard to beat in
practice. You grow a collection of decision trees, each on its
own bootstrap sample of the training rows, and at every split
inside each tree you let the algorithm pick from a random
subset of the features rather than from all of them. At
prediction time the trees vote. The randomness in the two
sampling steps is what stops the forest from agreeing too
closely with itself --- the resulting ensemble keeps the
interpretability of individual trees but throws away most of
their variance.

I picked it for the NIDS pipeline because of three properties
that mattered to me. The first is calibrated probabilities,
which gave me a sensible threshold knob (\code{0.8} by default)
between flagging too much and flagging too little. The second
is robustness: the algorithm tolerates noisy or partially
irrelevant features without obvious degradation, and I knew
upfront that not all six of my features (Section
\ref{sec:design-features}) would be equally informative on
every packet. The third is the \code{feature\_importances\_}
attribute on the trained model. Surfacing it as a bar chart on
the ML Models page (Figure~\ref{fig:ml-models}) means an
analyst can see, at a glance, that the classifier weighted
\code{inter\_arrival\_ms} heavily for a particular alert ---
that's actionable; a bare probability score isn't.

The downside is that it needs labelled training data. For this
thesis the training set is synthesised from known patterns and
that works fine. In production deployment, keeping labels
fresh enough to retrain on becomes the dominant ongoing cost,
which is one of the reasons I wanted the second algorithm
below as a backup.

\subsection{Isolation Forest}
\label{sec:bg-iforest}

The second algorithm came from a 2008 paper by Liu, Ting and
Zhou~\cite{liu2008isolation}, and the intuition behind it is
charming once you see it. Anomalies, the authors argue, are by
definition both rare and different, so a random partitioning
scheme should be able to isolate them faster than it isolates
the rest of the data. The construction reflects that idea
directly. Grow a forest of binary trees, each built by picking
a random feature and a random split point between that
feature's minimum and maximum value. To score a new sample,
drop it through every tree and measure how deep it had to
travel before reaching a leaf. Anomalies travel less. The
shorter the average path, the more anomalous the sample is
considered to be.

The advantage over a Random Forest is one I keep coming back
to: no labels. You only need a clean sample of normal traffic,
which is much easier to collect than a labelled attack corpus.
The disadvantage shows up at the edges of the normal
distribution, where rare-but-legitimate behaviour starts
looking suspicious in exactly the way attacks do. That's the
exact failure mode that Sommer and Paxson called
out~\cite{sommer2010outside} in their critique of academic
ML-IDS, and it's the reason I shipped \emph{both} algorithms
behind the same interface rather than committing to one. An
operator who has labelled data gets the supervised Random
Forest. One who doesn't gets the unsupervised Isolation
Forest. The rest of the pipeline doesn't change.

\section{Hybrid Detection}
\label{sec:bg-hybrid}

When both engines run on the same packet stream, the dispatcher can
produce three different alert categories:

\begin{description}
  \item[Signature-only:] the rule engine matched but the ML model
    classified the packet as benign. This is the highest-precision
    category --- the rule guarantees the packet matches a known
    pattern, and the model's disagreement carries no penalty.
  \item[ML-only:] no rule matched but the model's score exceeded the
    alert threshold. This category catches patterns for which no
    signature exists, at the cost of less explanation. An analyst
    looking at an ml-only alert knows the model thought the packet
    was malicious, but has to dig into the feature vector to see why.
  \item[Hybrid:] both engines flagged the same packet. This is the
    most reliable category. The rule contributes precision (a human
    decided this pattern is malicious) and the model contributes
    corroboration (the statistical evidence supports the rule).
\end{description}

The categorisation is performed in the dispatcher's
\code{\_evaluate} method. Each candidate alert carries an explicit
\code{detection\_method} field that propagates all the way to the UI,
where the analyst sees a coloured chip indicating the source. This
removes the black-box critique that has historically been levelled at
ML-based IDS~\cite{sommer2010outside}: the alert always tells the user
which engine produced it, and a click on the alert detail page shows
both the rule that matched and the ML score on the same screen.

\section{Forensics}
\label{sec:bg-forensics}

The forensic dimension is what distinguishes operational SIEM tools
(Wazuh, Splunk, IBM QRadar) from detection-only engines. An alert by
itself is just a row in a log; what makes it actionable is the
context. The classic forensic workflow is: see the alert, retrieve
the packets that produced it, eyeball them, decide if the alert is a
true or false positive, take the appropriate downstream action.

In a typical NIDS deployment, the retrieval step is delegated to a
full-packet-capture system that operates in parallel with the
detection engine. Moloch / Arkime is the canonical open-source
example. The architecture is sound but it involves running two
separate stacks, correlating their data by timestamp and 5-tuple, and
training the analyst on two different consoles. \nidsname{} takes a
different position: every alert carries its own packet capture
inline.

The mechanism is straightforward. The dispatcher maintains a
fixed-size ring buffer (\code{deque(maxlen=200)}) of the most recent
packets seen. When an alert is produced, the dispatcher slices the
last 50 entries off the ring and serialises them to a JSON Lines file
named after the alert ID. The path is stored on the
\code{Alert} row and a separate \code{PcapRecord} row indexes it for
the UI. Download is served by an authenticated FastAPI endpoint that
streams the file with a \code{FileResponse}. Section
\ref{sec:design-forensics} walks through the implementation.

JSON Lines was chosen over the libpcap binary format for two
practical reasons. JSONL is human-readable, so an analyst can open
the file in any text editor or pipe it through \code{jq} for
filtering. And JSONL does not depend on libpcap or wpcap being
installed at the time the file is created, which simplifies the
Windows deployment story considerably.
